% Plant Disease VR Experience -- Beamer (Metropolis) slides
% Compile with XeLaTeX (Overleaf: Menu -> Compiler -> XeLaTeX).
% Images are read from ../Tools/out/previews/ ; on Overleaf, upload the four
% scene_*.png files next to this .tex and the \graphicspath still finds them.

\documentclass[aspectratio=169,11pt]{beamer}
\usetheme[progressbar=frametitle, numbering=counter]{metropolis}
\usepackage{booktabs}
\usepackage{graphicx}
\usepackage{hyperref}
\graphicspath{{../Tools/out/previews/}{./}}

% ---- Colours: Metropolis dark teal + plant green + orange alerts ----
\definecolor{mDarkTeal}{HTML}{23373B}
\definecolor{PlantGreen}{HTML}{3B8C5A}
\definecolor{PlantGreenText}{HTML}{2D7048}
\definecolor{AlertOrange}{HTML}{A6560C}
\setbeamercolor{progress bar}{fg=PlantGreen, bg=PlantGreen!20}
\setbeamercolor{title separator}{fg=PlantGreen}
\setbeamercolor{alerted text}{fg=AlertOrange}
\setbeamercolor{example text}{fg=PlantGreenText}
\setbeamercolor{block title}{bg=mDarkTeal!12, fg=mDarkTeal}
\setbeamercolor{block body}{bg=mDarkTeal!5}
\setbeamercolor{block title example}{bg=PlantGreen!22, fg=PlantGreenText}
\setbeamercolor{block body example}{bg=PlantGreen!10}
\setbeamercolor{block title alerted}{bg=orange!20, fg=AlertOrange}
\setbeamercolor{block body alerted}{bg=orange!8}
\metroset{block=fill}

\newcommand{\preview}[2][\linewidth]{\includegraphics[width=#1]{#2}}
\newcommand{\figcap}[1]{{\scriptsize\textcolor{PlantGreenText}{Figure:} #1}}

\title{Plant Disease VR Experience}
\subtitle{Exploring plants in a virtual farm and an illustrative disease scenario in a lab}
\author{Bhargav M}
\institute{B.Tech Project (BTP) \textperiodcentered{} Interactive VR prototype}
\date{September 2026}

\begin{document}

% 1 ------------------------------------------------------------------
\maketitle
\note{0:25 -- Introduce the project as an interactive VR prototype.}

% 2 ------------------------------------------------------------------
\begin{frame}{Problem and Motivation}
  \begin{columns}[T]
    \begin{column}{0.44\textwidth}
      \begin{exampleblock}{Why it matters}
        {\Huge\textbf{20--40\%}}\\[4pt]
        of global crop yields are lost to plant pests and diseases each year~[1]
      \end{exampleblock}
      \begin{itemize}
        \item Leaf symptoms and disease progression are hard to communicate with static pictures alone.
        \item \textbf{Objective:} let a learner inspect 3D plants and explore one visual recovery example in VR.
      \end{itemize}
    \end{column}
    \begin{column}{0.53\textwidth}
      \preview{scene_spawn.png}\\
      \figcap{Built farm scene (Unity editor preview)}
    \end{column}
  \end{columns}
\end{frame}

% 3 ------------------------------------------------------------------
\begin{frame}{Literature Review}
  \begin{columns}[T]
    \begin{column}{0.31\textwidth}
      \begin{block}{Mohanty et al. (2016) [2]}
        \textbf{Image-based disease recognition}\\[3pt]
        Recognition from plant images is possible, but results can change outside controlled image conditions.
      \end{block}
    \end{column}
    \begin{column}{0.31\textwidth}
      \begin{block}{Cheng (2024) [3]}
        \textbf{Immersive VR for plant learning}\\[3pt]
        Plant-learning research has explored immersive VR as a way to engage learners.
      \end{block}
    \end{column}
    \begin{column}{0.31\textwidth}
      \begin{exampleblock}{This project}
        \textbf{VR exploration first}\\[3pt]
        Explore plants and visualize one example disease journey. Automated classification is future work.
      \end{exampleblock}
    \end{column}
  \end{columns}
\end{frame}

% 4 ------------------------------------------------------------------
\begin{frame}{Quick Recap: Previous Evaluation}
  From watching 360$^\circ$ video to exploring a real-time Unity scene
  \vspace{1em}
  \begin{columns}[c]
    \begin{column}{0.43\textwidth}
      \begin{block}{Previous evaluation}
        \textbf{360$^\circ$ VR Movie Player}
        \begin{itemize}
          \item Plays VR movies
          \item Completed as a separate project
        \end{itemize}
      \end{block}
    \end{column}
    \begin{column}{0.08\textwidth}
      \centering{\Huge\textcolor{PlantGreen}{$\Rightarrow$}}
    \end{column}
    \begin{column}{0.43\textwidth}
      \begin{exampleblock}{Current evaluation}
        \textbf{Plant Disease VR Experience}
        \begin{itemize}
          \item A new, real-time Unity project
          \item Selectable 3D plants in a farm and a lab
        \end{itemize}
      \end{exampleblock}
    \end{column}
  \end{columns}
\end{frame}

% 5 ------------------------------------------------------------------
\begin{frame}{Methodology: The Learner's Journey}
  \begin{columns}[T]
    \begin{column}{0.24\textwidth}
      \preview{scene_spawn.png}\\[2pt]
      \textcolor{PlantGreenText}{\textbf{1 \textperiodcentered{} Farm}}\\
      {\small Start in the farm and inspect a labeled plant.}
    \end{column}
    \begin{column}{0.24\textwidth}
      \preview{scene_specimens.png}\\[2pt]
      \textcolor{PlantGreenText}{\textbf{2 \textperiodcentered{} Select plant}}\\
      {\small Select a plant and enter the lab through a world-space button.}
    \end{column}
    \begin{column}{0.24\textwidth}
      \preview{scene_lab_plant.png}\\[2pt]
      \textcolor{PlantGreenText}{\textbf{3 \textperiodcentered{} Lab}}\\
      {\small Inspect its 3D model and information panel.}
    \end{column}
    \begin{column}{0.24\textwidth}
      \preview{scene_lab_left.png}\\[2pt]
      \textcolor{PlantGreenText}{\textbf{4 \textperiodcentered{} Tomato demo}}\\
      {\small Six-week Early Blight visualization, then return to the farm.}
    \end{column}
  \end{columns}
  \vspace{0.8em}
  \begin{alertblock}{Design boundary}
    The six-week values are illustrative sample data. The experience does not infer disease from a camera image.
  \end{alertblock}
\end{frame}

% 6 ------------------------------------------------------------------
\begin{frame}{Implementation: Unity and XR Architecture}
  \begin{columns}[T]
    \begin{column}{0.42\textwidth}
      \begin{description}[Transitions]
        \item[Engine] Unity 6000.5.5f1, Universal Render Pipeline
        \item[XR] OpenXR, XR Interaction Toolkit, XR Hands
        \item[Scenes] One persistent XR bootstrap; farm and lab load as separate environment scenes
        \item[Transitions] Fade, keeping the selected plant
      \end{description}
    \end{column}
    \begin{column}{0.55\textwidth}
      \centering
      \colorbox{mDarkTeal}{\parbox{0.9\linewidth}{\centering\color{white}\textbf{XRBootstrap}\\{\scriptsize persistent XR rig \textperiodcentered{} camera \textperiodcentered{} interaction}}}\\[2pt]
      {\color{PlantGreen}$\Downarrow$\hspace{0.4\linewidth}$\Downarrow$}\\[2pt]
      \begin{minipage}{0.48\linewidth}\centering
        \preview{scene_spawn.png}\\{\scriptsize\textbf{Farm} (Envirornment.unity)}
      \end{minipage}\hfill
      \begin{minipage}{0.48\linewidth}\centering
        \preview{scene_lab_plant.png}\\{\scriptsize\textbf{Lab}}
      \end{minipage}\\[3pt]
      \figcap{Unity editor previews}
    \end{column}
  \end{columns}
\end{frame}

% 7 ------------------------------------------------------------------
\begin{frame}{Implementation: Plants and Asset Workflow}
  \begin{columns}[T]
    \begin{column}{0.5\textwidth}
      \preview{scene_specimens.png}\\
      \figcap{Farm scene (Unity editor preview)}
    \end{column}
    \begin{column}{0.47\textwidth}
      \begin{exampleblock}{Mesh reduction (Tomato, farm)}
        {\LARGE\textbf{1.62M $\rightarrow$ 8,000}} triangles
      \end{exampleblock}
      \begin{enumerate}\small
        \item Supplied FBX / GLB plant models
        \item Heavy meshes reduced in Blender
        \item Textures and URP materials in Unity
        \item Farm and lab prefabs; billboards at a distance
      \end{enumerate}
    \end{column}
  \end{columns}
  \vspace{0.6em}
  {\small\textcolor{PlantGreenText}{\textbf{Integrated plants:}} Apple \textperiodcentered{} Bell Pepper \textperiodcentered{} Cherry \textperiodcentered{} Corn (Maize) \textperiodcentered{} Potato \textperiodcentered{} Strawberry \textperiodcentered{} Tomato}
\end{frame}

% 8 ------------------------------------------------------------------
\begin{frame}{Implementation: Lab and Tomato Demo}
  \begin{columns}[T]
    \begin{column}{0.48\textwidth}
      \preview{scene_lab_plant.png}\\
      \figcap{Central plant pedestal (Unity editor preview)}
    \end{column}
    \begin{column}{0.48\textwidth}
      \preview{scene_lab_left.png}\\
      \figcap{Selection panel and week timeline (Unity editor preview)}
    \end{column}
  \end{columns}
  \vspace{0.5em}
  \begin{itemize}\small
    \item World-space selection, information and recovery panels around a central pedestal
    \item Six selectable weeks of Early Blight; health bars and plant appearance change by week
    \item Week 3 is labeled \textbf{Demo treatment applied}
  \end{itemize}
  {\small\alert{\textit{On screen: ``Illustrative prototype. Not diagnostic or treatment advice.''}}}
\end{frame}

% 9 ------------------------------------------------------------------
\begin{frame}{Results: Verified Project Outputs}
  \begin{columns}[T]
    \begin{column}{0.62\textwidth}
      \small
      \begin{tabular}{@{}p{0.44\linewidth}p{0.5\linewidth}@{}}
        \toprule
        \textbf{Verified in the project folder} & \textbf{Evidence} \\
        \midrule
        Three enabled scenes & XR bootstrap, farm and lab in Build Settings \\
        Seven integrated crops & Plant catalog plus 14 farm/lab prefabs \\
        Tomato six-week example & Recovery data asset and UI code \\
        \textbf{16/16 EditMode tests passed} & Unity Editor test log \\
        Quest APK and Windows .exe & Builds/Quest and Builds/Windows \\
        \bottomrule
      \end{tabular}
    \end{column}
    \begin{column}{0.34\textwidth}
      \begin{exampleblock}{Unit tests}
        {\Huge\textbf{16/16}}\\[3pt]
        {\small EditMode tests passed: catalog, selection and timeline logic}
      \end{exampleblock}
    \end{column}
  \end{columns}
\end{frame}

% 10 -----------------------------------------------------------------
\begin{frame}{Analysis: Limits of the Current Evaluation}
  \begin{columns}[T]
    \begin{column}{0.47\textwidth}
      \begin{exampleblock}{Verified now}
        \begin{itemize}\small
          \item Three scenes, seven crops and the Tomato example exist
          \item 16 of 16 EditMode tests pass
          \item Quest APK and Windows executable are packaged
        \end{itemize}
      \end{exampleblock}
    \end{column}
    \begin{column}{0.47\textwidth}
      \begin{alertblock}{Next validation}
        \begin{itemize}\small
          \item No recorded Quest 2/3 or PC headset performance results yet
          \item Controller and hand interaction need on-device checks
          \item Grape and Peach models not yet in the selector
          \item Tomato percentages are sample values; no classification or treatment guidance
        \end{itemize}
      \end{alertblock}
    \end{column}
  \end{columns}
  \vspace{1em}
  \centering\textit{A functional Unity prototype, not yet a validated plant diagnosis system.}
\end{frame}

% 11 -----------------------------------------------------------------
\begin{frame}{Conclusion and Future Work}
  \begin{exampleblock}{Current result}
    A farm-to-lab VR prototype with seven selectable plants and one illustrative Tomato example.
  \end{exampleblock}
  \begin{columns}[T]
    \begin{column}{0.24\textwidth}
      \begin{block}{Next validation}\scriptsize Quest 2/3 and PC VR: frame rate, interaction, learner usability\end{block}
    \end{column}
    \begin{column}{0.24\textwidth}
      \begin{block}{Next plants}\scriptsize Optimize and integrate Grape and Peach\end{block}
    \end{column}
    \begin{column}{0.24\textwidth}
      \begin{block}{Disease system}\scriptsize Five verified diseases per plant, plus a healthy class\end{block}
    \end{column}
    \begin{column}{0.24\textwidth}
      \begin{block}{Interactive care}\scriptsize Choose treatments and, where appropriate, fertilizer type and proportion; see a simulated response\end{block}
    \end{column}
  \end{columns}
\end{frame}

% 12 -----------------------------------------------------------------
\begin{frame}[allowframebreaks]{References}
  \small
  \begin{itemize}
    \item[{[1]}] FAO. \emph{Understanding the context: Pest and Pesticide Management.} \url{https://www.fao.org/pest-and-pesticide-management/about/understanding-the-context/en/}
    \item[{[2]}] S.~P. Mohanty, D.~P. Hughes and M.~Salath\'e. Using deep learning for image-based plant disease detection. \emph{Frontiers in Plant Science}, 7:1419, 2016.
    \item[{[3]}] K.-H. Cheng. Development of an immersive virtual reality system for learning about plants in primary education. \emph{Educational Technology Research and Development}, 72:845--867, 2024.
    \item[{[4]}] Unity project source, assets, editor test log and build outputs, reviewed 25 September 2026.
  \end{itemize}
\end{frame}

% 13 -----------------------------------------------------------------
\begin{frame}[standout]
  Thank you\\[0.5em]
  {\normalsize Questions?}
\end{frame}

\end{document}
